\chapter*{符号索引}
\addcontentsline{toc}{chapter}{符号索引}

下表按对象类型汇总各章使用的符号、单位和记号约定。数学定义与物理意义以正文首次出现处为准；单位制统一为 SI，并显式保留 $c$、$\hbar$、$\varepsilon_0$、$\mu_0$ 与 $k_B$。

\section*{时空、规范与傅里叶变换约定}
\begin{longtable}{p{0.23\textwidth}p{0.53\textwidth}p{0.16\textwidth}}
\toprule
符号 & 数学定义与物理意义 & SI 量纲\\
\midrule
\endfirsthead
\toprule
符号 & 数学定义与物理意义 & SI 量纲\\
\midrule
\endhead
$x^\mu$ & $x^\mu=(ct,\mathbf r)$，Minkowski 四维坐标。 & $\mathrm m$\\
$\eta_{\mu\nu}$ & $\eta_{\mu\nu}=\operatorname{diag}(1,-1,-1,-1)$，全文唯一 Minkowski 度规。 & 1\\
$\partial_\mu,\partial^\mu$ & $\partial_\mu=(c^{-1}\partial_t,\nabla)$、$\partial^\mu=(c^{-1}\partial_t,-\nabla)$，由度规升降指标。 & $\mathrm{m^{-1}}$\\
$p^\mu$ & $p^\mu=(E/c,\mathbf p)$，任意有质量或无质量粒子的物理四动量；其平方 $p^2=\eta_{\mu\nu}p^\mu p^\nu$。 & kg m s$^{-1}$\\
$\kappa^\mu$ & $\kappa^\mu=(\omega/c,\boldsymbol\kappa)$，傅里叶四波矢；真实光子的物理四动量统一写 $k^\mu=\hbar\kappa^\mu$。 & m$^{-1}$\\
$\mathbf G,G_m$ & 晶格理论专用的倒格矢与一维倒格波数；它们只表示离散平移对称性产生的倒格对象，电子相应交换机械动量为 $\hbar\mathbf G$ 或 $\hbar G_m$。 & m$^{-1}$\\
$e,\;q_{\mathrm e}$ & $e>0$ 为基本电荷的正值；电子电荷固定为 $q_{\mathrm e}=-e$。 & C\\
$A^\mu,A_\mu$ & $A^\mu=(\phi/c,\mathbf A)$、$A_\mu=(\phi/c,-\mathbf A)$，电磁四势。各分量具有同一 SI 量纲。 & $\mathrm{V\,s\,m^{-1}}$\\
$F^{\mu\nu}$ & $F^{\mu\nu}=\partial^\mu A^\nu-\partial^\nu A^\mu$；$F^{0i}=-E_i/c$、$F^{ij}=-\epsilon_{ijk}B_k$。 & T\\
${}^\star\!F^{\mu\nu}$ & 电磁场强的 Hodge 对偶，${}^\star\!F^{\mu\nu}=\epsilon^{\mu\nu\rho\sigma}F_{\rho\sigma}/2$；星号只表示 Hodge 对偶，波浪号只用于 Fourier 变换。 & T\\
$D_\mu$ & $D_\mu=\partial_\mu+\ii q_{\mathrm e} A_\mu/\hbar$，电子的最小耦合协变导数。 & $\mathrm{m^{-1}}$\\
$\chi(x)$ & 规范函数；全文采用 $A_\mu' = A_\mu-\partial_\mu\chi$ 与 $\psi'=\ee^{\ii q_{\mathrm e}\chi/\hbar}\psi$。 & $\mathrm{V\,s}$\\
$\widetilde f(\omega)$ & 时间傅里叶约定：$\widetilde f(\omega)=\int\dd t\,f(t)\ee^{+\ii\omega t}$，$f(t)=\int \dd\omega\,\widetilde f(\omega)\ee^{-\ii\omega t}/(2\pi)$。 & 由 $f$ 决定\\
$\widetilde f_z(\kappa)$ & 一维空间傅里叶约定：$\widetilde f_z(\kappa)=\int\dd\zeta\,f_z(\zeta)\ee^{-\ii\kappa\zeta}$，$f_z(\zeta)=\int\dd\kappa\,\widetilde f_z(\kappa)\ee^{+\ii\kappa\zeta}/(2\pi)$；除晶格理论专用的倒格矢 $\mathbf G$ 或一维倒格波数 $G_m$ 外，所有具有 $\mathrm{m^{-1}}$ 量纲的空间波数统一使用 $\kappa$ 字母族，三维反变换相应使用 $(2\pi)^{-3}$。 & 由 $f_z$ 决定\\
\bottomrule
\end{longtable}

高频字母按对象类型保留：一般具有 $\mathrm{m^{-1}}$ 量纲的波矢/波数使用 $\kappa$ 字母族，共振线宽与衰减率使用大写 $\Gamma$ 字母族，统计 cumulant 使用 $\mathcal C_n$，高斯环境二点函数的对称/反对称核使用 $C_{\rm sym},C_{\rm asym}$；倒格矢作为晶格专用对象保留 $\mathbf G$（一维写 $G_m$）。自由电子量子光学部分的近似控制量统一写成带物理下标的 $\epsilon$，探测或收集概率使用 $\mathcal Q^{\rm det}$、$\mathcal Q_{\rm col}$，通道分支比使用 $\mathcal B$；$\eta_{\mu\nu}$ 保留给 Minkowski 度规。$p^\mu,k^\mu,q^\mu$ 都是物理动量量纲的四矢量，具体粒子与传递方向在每个过程中定义；对真实光子有 $k^\mu=\hbar\kappa^\mu$。能量四矢量另以 $P_E=cp$ 等显式关系定义。无量纲光学、传播或格点累积相位使用 $\varphi$ 字母族；电磁标势保持为 $\phi$。规范耦合中，$g_{\rm G}$ 只表示尚未指定具体群时的一般规范耦合；$g$、$g_s$ 分别保留给标准模型弱 $SU(2)_L$ 与 QCD；重整化群的一般单耦合示例使用 $g_{\rm R}$，QPINEM 单程交换耦合固定为 $g_{\mathrm Q}$。速度比写 $\beta_{\rm rel}$，而经典 PINEM 累积耦合写 $\beta$。非阿贝尔规范联络统一写 $\mathscr C_\mu$，其曲率写 $\mathcal F_{\mu\nu}$；花体 $\mathcal A$ 只用于跃迁或散射振幅。几何/脉冲面积若需要单独记号则使用手写体 $\mathscr A$，避免与振幅混用。

\section*{自由电子与精确运动学}
\begin{longtable}{p{0.23\textwidth}p{0.53\textwidth}p{0.16\textwidth}}
\toprule
符号 & 数学定义与物理意义 & SI 量纲\\
\midrule
\endfirsthead
\toprule
符号 & 数学定义与物理意义 & SI 量纲\\
\midrule
\endhead
$m_{\mathrm e}$ & 电子静质量。 & kg\\
$E(\mathbf p)$ & $E(\mathbf p)=\sqrt{m_{\mathrm e}^2c^4+c^2\mathbf p^2}$，正能自由电子的精确相对论能量。 & J\\
$\mathbf p_0,E_0$ & 入射波包中心动量与对应总能量 $E_0=E(\mathbf p_0)$。 & kg m s$^{-1}$；J\\
$\gamma_{\rm L}$ & 同一 Lorentz 因子可写为 $\gamma_{\rm L}(\mathbf p)=E(\mathbf p)/(m_{\mathrm e} c^2)$ 或 $\gamma_{\rm L}(v)=(1-v^2/c^2)^{-1/2}$。 & 1\\
$\gamma_0$ & $\gamma_0=E_0/(m_{\mathrm e} c^2)$，中心 Lorentz 因子，即 $\gamma_{\rm L}(\mathbf p_0)$。 & 1\\
$\beta_{\rm rel}$ & $\beta_{\rm rel}=v/c$，只表示速度与光速之比；经典 PINEM 的无量纲耦合始终写 $\beta$，二者不混用。 & 1\\
$q^\mu$ & $q^\mu=p_f^\mu-p_i^\mu$，精确四动量转移；电子电荷始终写 $q_{\mathrm e}$，不以 $q$ 表示。 & kg m s$^{-1}$\\
$\mathcal Q$ & $\mathcal Q=\sqrt{-q^2}>0$，类空硬动量的大小；其能量形式唯一写 $E_{\mathcal Q}=c\mathcal Q$。 & kg m s$^{-1}$\\
$\mathbf v_0$ & $\mathbf v_0=\nabla_{\mathbf p}E|_{\mathbf p_0}=c^2\mathbf p_0/E_0$，中心群速度。 & m s$^{-1}$\\
$m_\perp$ & $m_\perp=\gamma_0m_{\mathrm e}$，相对论色散横向二阶曲率对应的有效质量。 & kg\\
$m_\parallel$ & $m_\parallel=\gamma_0^3m_{\mathrm e}$，相对论色散纵向二阶曲率对应的有效质量。 & kg\\
$\hat{\mathcal V},\hat{\mathcal O}$ & Foldy--Wouthuysen 分块中的偶算符与奇算符；静态最小耦合时 $\hat{\mathcal V}=q_{\mathrm e}\phi$、$\hat{\mathcal O}=c\boldsymbol\alpha\cdot\hat{\boldsymbol\Pi}$。 & J\\
$\kappa_{\rm e0}$ & $\kappa_{\rm e0}=p_0/\hbar$，正能窄带电子包络的中心载波波数。 & m$^{-1}$\\
$f_{\mathrm e}$ & 正能窄带投影后缓变的电子纵向包络；它不表示规范耦合或 QPINEM 交换强度。 & 依归一化而定\\
$\lambda_{\rm dB}$ & $\lambda_{\rm dB}=h/p_0$，中心 de Broglie 波长。 & m\\
$|\mathbf p,s,+\rangle$ & 正能自由 Dirac 动量--自旋本征态；连续态采用相应定义式给出的动量 $\delta$ 函数归一化。 & 广义态\\
\bottomrule
\end{longtable}

\section*{场、介质、模式与联合动力学}
\begin{longtable}{p{0.23\textwidth}p{0.53\textwidth}p{0.16\textwidth}}
\toprule
符号 & 数学定义与物理意义 & SI 量纲\\
\midrule
\endfirsthead
\toprule
符号 & 数学定义与物理意义 & SI 量纲\\
\midrule
\endhead
$\mathsf L_{\rm EM}(\omega)$ & 一般线性 Maxwell 算符；材料响应、几何与边界条件均进入该算符，其推迟逆算符定义电磁 Green 张量。 & 依表示而定\\
$\mathbf G^R(\mathbf r,\mathbf r';\omega)$ & 推迟并矢 Green 张量，即 $\mathsf L_{\rm EM}(\omega)$ 的推迟逆核；$\operatorname{Im}\mathbf G^R=[\mathbf G^R-(\mathbf G^R)^\dagger]/(2\ii)$，其中核的伴随同时交换两端坐标并作转置复共轭。 & 依 Green 定义\\
$\mathcal M$ & 约化不变量散射振幅；$2\to2$ QFT 截面采用无量纲 $\mathcal M$ 约定，所有 SI 面积因子显式留在相空间与通量前因子中。 & 1\\
$\mathcal A_{f\leftarrow i}$ & 一般复跃迁/散射振幅族；具体下标说明初末态或通道，例如 $\mathcal A_{\ell|n}$。其量纲随所用连续态或离散态归一化而定；QFT 的约化不变量振幅专写 $\mathcal M$。 & 依归一化而定\\
$\dd\mathfrak P_n$ & $n$ 体 Lorentz 不变末态相空间测度；专用于散射与衰变相空间，不与 Higgs 二重态 $\Phi$ 复用。 & 按粒子数由定义式决定\\
$\varphi_{\rm ch},m_{\rm ch}$ & QFT 基础与标量玩具模型中的一般复标量物质场及其静质量；大写 $\Phi$ 专留给标准模型 Higgs 二重态。 & $\sqrt{\mathrm{J/m}}$；kg\\
$\psi_{\rm 1p}(\mathbf r)$ & 固定一粒子扇区中的普通复数值 Schrödinger 波函数；与场算符 $\hat\Psi(\mathbf r)$ 区分。 & m$^{-3/2}$（归一化到 1 时）\\
$\varphi_{\rm ssb},g_{\rm G}$ & 一般自发对称性破缺模型中的复标量场与无量纲规范耦合；只用于通用 Goldstone/Higgs 机制，不与标准模型 $\Phi,g$ 复用。 & $\sqrt{\mathrm{J/m}}$；1\\
$g_{\rm R}$ & 重整化群的一般单耦合示例中使用的重整化无量纲耦合，$\beta_{g_{\rm R}}=\mu\,\dd g_{\rm R}/\dd\mu|_{\rm bare}$；它不表示任何特定规范群的规范耦合。 & 1\\
$\Gamma^\mu(p',p)$ & 剥去电荷因子的精确单粒子不可约电磁顶角；树级为 $\gamma^\mu$，并满足 Ward--Takahashi 恒等式。 & 1\\
$T^a,f^{abc}$ & $SU(N)$ 基本表示生成元与结构常数；$\Tr(T^aT^b)=\delta^{ab}/2$、$[T^a,T^b]=\ii f^{abc}T^c$。 & 1\\
$g_s,\alpha_s$ & QCD 无量纲规范耦合与 $\alpha_s=g_s^2/(4\pi)$。 & 1\\
$\psi_f(x)$ & 味 $f$ 的 QCD 夸克 Dirac 场；$q^\mu$ 只保留四动量转移含义。 & m$^{-3/2}$\\
$\mathscr C_\mu^a$ & QCD 几何归一化的色规范联络；与电磁 SI 四势 $A_\mu$ 区分。 & m$^{-1}$\\
$\mathcal F_{\mu\nu}^a$ & QCD 色场强，由 $[D_\mu,D_\nu]=\ii g_s\mathcal F_{\mu\nu}$ 定义；与电磁场强 $F_{\mu\nu}$ 区分。 & m$^{-2}$\\
${}^\star\!\mathcal F_{\mu\nu}^a$ & QCD 色场强的 Hodge 对偶，${}^\star\!\mathcal F^{a\mu\nu}=\epsilon^{\mu\nu\rho\sigma}\mathcal F_{\rho\sigma}^a/2$；星号表示对偶，波浪号只保留给 Fourier 变换。 & m$^{-2}$\\
$g,g',g_{\rm em}$ & $SU(2)_L$、$U(1)_Y$ 与未破缺 $U(1)_{\rm em}$ 的无量纲规范耦合；$g_{\rm em}=g\sin\theta_W=g'\cos\theta_W$。 & 1\\
$A_\mu^{\rm geo},F_{\mu\nu}^{\rm geo}$ & 标准模型电弱质量基中的几何归一化光子联络及其场强；$A_\mu^{\rm geo}=\sqrt{\varepsilon_0c/\hbar}\,A_\mu$、$F_{\mu\nu}^{\rm geo}=\sqrt{\varepsilon_0c/\hbar}\,F_{\mu\nu}$，与前文 SI 电磁四势、场强严格区分。 & m$^{-1}$；m$^{-2}$\\
$T^I,Y,Q$ & 弱同位旋生成元、超荷与电荷生成元；全文采用 $Q=T^3+Y$。 & 1\\
$\Phi$ & 标准模型 Higgs 二重态，规范表示 $(\mathbf1,\mathbf2)_{1/2}$。 & $\sqrt{\mathrm{J/m}}$\\
$v,v_E$ & Higgs 真空振幅及其能量形式；$v_E=\sqrt{\hbar c}\,v$。 & $\sqrt{\mathrm{J/m}}$；J\\
$\lambda_H,\mu_H^2$ & Higgs 四次耦合与二次参数，$V=-\mu_H^2\Phi^\dagger\Phi+\lambda_H(\Phi^\dagger\Phi)^2/(\hbar c)$。 & 1；m$^{-2}$\\
$\mu_E^2$ & Higgs 二次参数的能量平方形式，$\mu_E^2=(\hbar c)^2\mu_H^2$；只表示 Higgs 势参数。 & J$^2$\\
$E_{\rm R}$ & 重整化能量尺度，唯一地定义为 $E_{\rm R}=c\mu_{\rm DR}$；不与 Higgs 的 $\mu_E^2$ 混用。 & J\\
$\mu_{\rm F},E_{\rm F}$ & 因子化动量与能量尺度，$E_{\rm F}=c\mu_{\rm F}$；与重整化尺度 $\mu_{\rm DR}$ 分开。 & kg m s$^{-1}$；J\\
$\Lambda_*,\mu_{\rm M},\Lambda_E$ & 重自由度有效场论的逆长度尺度、匹配动量与匹配能量尺度；$\mu_{\rm M}=\hbar\Lambda_*$，$\Lambda_E=c\mu_{\rm M}=\hbar c\Lambda_*$。 & m$^{-1}$；kg m s$^{-1}$；J\\
$\varphi_E$ & 常数中性 Higgs 背景的能量变量，$\varphi_E=\sqrt{\hbar c}\,\varphi$。 & J\\
$\mathbf E(\mathbf r,t)$ & 时间域实电场。全文凡未带频率下标的 $\mathbf E$ 均指这一对象。 & V m$^{-1}$\\
$\mathbf E_\omega(\mathbf r)$ & 单频正频分量的复空间幅度，由 $\mathbf E^{(+)}(\mathbf r,t)=\mathbf E_\omega(\mathbf r)\ee^{-\ii\omega t}$ 定义；不是另一个独立的场变量。 & V m$^{-1}$\\
$E_{\omega,z}(\mathbf r)$ & $\mathbf E_\omega$ 沿电子参考轨迹纵向轴的分量；PINEM 同步积分统一使用此记号。 & V m$^{-1}$\\
$\widetilde{\mathbf E}(\mathbf r,\omega)$ & 对一般非单色时域电场按全书 Fourier convention 作时间变换得到的频谱幅；与单频 phasor $\mathbf E_\omega$ 区分。 & V s m$^{-1}$\\
$\widetilde{\mathbf J}(\mathbf r,\omega)$ & 时域电流密度 $\mathbf J(\mathbf r,t)$ 的时间 Fourier 变换；在 EELS/辐射功率谱中使用。 & C m$^{-2}$\\
$\omega_\lambda$ & 第 $\lambda$ 个玻色正常模或环境库分量的正角频率。 & s$^{-1}$\\
$\hat a_\lambda,\hat a_\lambda^\dagger$ & 第 $\lambda$ 个玻色模的湮灭与产生算符。 & 1\\
$\mathcal G_{\mathbf p'\mathbf p}^{(\lambda)}$ & Dirac-$\delta$ 归一化连续电子态上的一量子耦合分布核；只有与 $\dd^3 p$ 测度组合后才有普通速率量纲。 & $[\mathcal G\,\dd^3 p]={\rm s^{-1}}$\\
$\mathcal G_{fi}^{(\lambda)},\mathcal G_\ell$ & 有限盒或正交波包/离散通道归一化后的普通一量子耦合率矩阵元。 & s$^{-1}$\\
$J(\omega)$ & 在给定模式测度和耦合约定下定义的环境谱密度。 & 依约定而定\\
$P,Q$ & $P+Q=I$、$PQ=0$ 的 Feshbach 目标/被消去投影算符。 & 1\\
$H_{\rm eff}(z)$ & $PHP+PHQ(z-QHQ)^{-1}QHP$，精确的能量依赖 Feshbach--Schur 有效哈密顿量。 & J\\
\bottomrule
\end{longtable}

\section*{离散通道、受控极限与测量}
\begin{longtable}{p{0.23\textwidth}p{0.53\textwidth}p{0.16\textwidth}}
\toprule
符号 & 数学定义与物理意义 & SI 量纲\\
\midrule
\endfirsthead
\toprule
符号 & 数学定义与物理意义 & SI 量纲\\
\midrule
\endhead
$|\ell\rangle$ & 谱投影后以有限盒或等价正交波包归一化的第 $\ell$ 个电子通道；不是连续分布态的“点值”，也不是电子内部能级。 & 归一化态\\
$\mathcal E_\ell$ & $\mathcal E_\ell=E(\mathbf p_\ell)$，真实相对论通道能量；始终与理想等间隔能量梯区分。 & J\\
$\Delta_\ell$ & $\Delta_\ell=(\mathcal E_{\ell+1}-\mathcal E_\ell)/\hbar-\omega$，真实逐边失谐。 & s$^{-1}$\\
$E_\ell^{(0)}$ & $E_0+\ell\hbar\omega$，仅在无反冲、均匀梯有效模型中使用的理想能量。 & J\\
$\hat b,\hat b^\dagger$ & 只有在正能谱边界不可见且平移对称条件成立后，才使用的双边电子通道平移算符。 & 1\\
$g_{\mathrm Q}$ & 单模量子交换在有限作用窗口内积分得到的无量纲单程耦合；不与耦合速率 $\mathcal G$ 混用。 & 1\\
$\mathscr U_{\mu\nu}^{s_f s_i}(\mathbf p_f,\mathbf p_i;T)$ & 正能一电子扇区与任意辅助谱之间的有限时间联合通道核；离散与连续辅助谱统一由谱测度 $\dd\mu_X$ 处理。 & 分布核，随连续基归一化而定\\
$\delta_X(\mu,\nu)$ & 辅助系统谱测度 $\dd\mu_X$ 上的恒等核，满足 $\int\dd\mu_X(\mu)\delta_X(\mu,\nu)f(\mu)=f(\nu)$。 & 谱测度的逆量纲\\
$\beta$ & 经典相干场支路的无量纲同步耦合；受控经典缩放下 $\beta=\alpha_{\rm coh}g_{\mathrm Q}$。 & 1\\
$\Gamma_c$ & 开放共振或 QNM 的角频率线宽，复极点写为 $\widetilde\omega_c=\omega_c-\ii\Gamma_c/2$；不使用 $\kappa$ 表示线宽。 & s$^{-1}$\\
$\Gamma_{\rm mat},\Gamma_{\rm ph},\Gamma_D$ & 材料、光学声子与 Drude 阻尼率；所有耗散率统一使用大写 $\Gamma$ 家族。 & s$^{-1}$\\
$\rho_{ij}$ & 密度矩阵元；$i\neq j$ 时表示基底态之间的相干，不另用 $\gamma$ 表示相干度。 & 1\\
$\mathcal C_n^{(E)}$ & 电子能量改变的第 $n$ 阶 cumulant；离散边带极限满足 $\mathcal C_n^{(E)}=(\hbar\omega)^n\mathcal C_n^{(\ell)}$。 & J$^n$\\
$C_{\rm sym},C_{\rm asym}$ & 高斯环境二点函数的对称涨落核与反对称响应核；分别由反对易子与对易子定义，并贯穿影响泛函、FDT 与 Born 记忆方程。 & 由耦合环境算符决定\\
$\mathscr A_{\rm hop}$ & $\mathscr A_{\rm hop}=|\Omega|T/2$，常耦合有限反冲格点的无量纲跳跃面积；平移极限中等于 $|\beta|$。 & 1\\
$\kappa_z,\;\omega_{\rm r}$ & $\kappa_z$ 是单频纵向交换波数，单步动量为 $\hbar\kappa_z$；$\omega_{\rm r}=\hbar\kappa_z^2/(2m_\parallel)$ 是由相对论纵向色散曲率产生的反冲频率。 & m$^{-1}$；s$^{-1}$\\
$\kappa_{\rm L},\kappa_{\rm KD},\omega_{\rm KD}$ & Kapitza--Dirac 驻波中单束激光波数、强度光栅波数与单光子反冲频率；$\kappa_{\rm KD}=2\kappa_{\rm L}$，$\omega_{\rm KD}=\hbar\kappa_{\rm L}^2/(2m_{\mathrm e})$。 & m$^{-1}$；m$^{-1}$；s$^{-1}$\\
$r_{\rm ph}$ & 弱交换 QPINEM 中的正整数光子关联/电子净增益阶数；只作阶数指标，不表示空间半径。 & 1\\
$\varphi_{\rm lin},\varphi_{\rm r}$ & $\varphi_{\rm lin}=\Delta_{\rm lin}T$、$\varphi_{\rm r}=\omega_{\rm r}T$，分别为线性失谐与反冲曲率在作用时间内累积的无量纲相位。 & 1\\
$B_{\Delta\ell}(L)$ & 传播距离 $L$ 处电子纵向密度的第 $\Delta\ell$ 阶傅里叶聚束因子，其中 $\Delta\ell=\ell-\ell\prime\in\mathbb Z$ 是能量格指标差；$B_{-\Delta\ell}=B_{\Delta\ell}^*$，同时等于相应副对角线的相干和。 & 1\\
$T_T,L_T$ & $T_T=2\pi/\omega_{\rm r}$、$L_T=v_0T_T$，二阶反冲色散下电子能量梳的 Talbot 复现时间与长度。 & s；m\\
$\epsilon_{\rm Qrec}^{(T)},\epsilon_{\rm Qrec}^{(\Omega)}$ & 有限反冲 QPINEM 中分别比较相邻交换边失谐与作用时间、局域交换速率的无量纲控制量。 & 1\\
$R_E(\Delta E)$ & 电子谱仪能量响应函数，按 $\int \dd E\,R_E(E)=1$ 归一化，用于把理想输出谱卷积为实验读出概率。 & J$^{-1}$\\
$\mathcal B_{\rm rad}(\omega)$ & 同一总 Green 张量与源问题中远场辐射概率占总电子能量损失概率的频率分辨分支比。 & 1\\
$\mathcal Q_\ell^{\rm det},\mathcal Q_E^{\rm det},\mathcal Q_\gamma^{\rm det},\mathcal Q_{\rm col}$ & 分别表示离散电子通道、连续电子能量、光子散射末态与收集子空间的无量纲探测/保留概率；效率统一使用 $\mathcal Q$ 字母族。 & 1\\
$\hat\rho$ & 量子密度算符；电荷密度另写 $\hat\rho_q(\mathbf r)$，不与概率密度混用。 & 1\\
$\Pi_x^{(S)}$ & 系统侧测量结果 $x$ 的正算符值测度（POVM）效应，满足 $\Pi_x^{(S)}\succeq0$ 与 $\sum_x\Pi_x^{(S)}=I_S$；测量仪侧效应写 $\Pi_x^{(M)}$。 & 1\\
$\mathcal I_x$ & 量子仪器中按结果分辨的完全正映射；既决定 $p(x)$ 又决定测后态。 & 超算符\\
\bottomrule
\end{longtable}

下表把常用近似的控制量、保留层级和失效时应恢复的母方程并列，便于直接比较不同约化的适用范围。

\begin{longtable}{@{}p{0.16\textwidth}p{0.32\textwidth}p{0.19\textwidth}p{0.20\textwidth}@{}}
\toprule
约化 & 控制量/尺度关系 & 保留层级 & 条件失效时的母方程\\
\midrule
\endfirsthead
\toprule
约化 & 控制量/尺度关系 & 保留层级 & 条件失效时的母方程\\
\midrule
\endhead
Schrieffer--Wolff（SW）低维化 & $\epsilon_{\rm SW}=\|PHQ\|/\Delta_Q\ll1$ & 给定阶的分块对角化 & 精确 Feshbach 预解式\\
单模 & $\epsilon_{\rm sm}(T)\ll1$，且 $\Delta\omega_{\rm sep}$ 大于线宽/耦合带宽 & 孤立模式 & 连续多模联合动力学\\
旋波近似 & $\epsilon_{\rm RWA}\ll1$，并以 $\epsilon_{\rm RWA}^{\rm dyn}\ll1$ 控制有限窗快相位积分与二阶虚反旋修正 & 共振频率块 & 完整四项含时耦合\\
程函 & $\epsilon_{\rm eik}=\max(\epsilon_u,\epsilon_\perp,\epsilon_p,\epsilon_{\rm WKB})\ll1$，并以 $\epsilon_{\rm eik}^{\rm op}$ 控制累计传播误差 & 固定轨迹相位输运 & 精确 Dirac 相互作用绘景传播子\\
边带概率解释 & $\epsilon_{\rm ov}\ll1$、$\epsilon_{\rm rec,eik}\ll1$、$\epsilon_{\rm sb,tail}\ll1$ & 周期相位的 Fourier 系数可解释为可分辨能量边带概率 & 连续动量卷积与真实 $E(p)$\\
无反冲均匀梯 & $\epsilon_{\rm tr}=\max(\epsilon_{\mathcal G},\epsilon_\Delta)\ll1$；量子支路另需 $\epsilon_{\rm tr}^{\rm dyn}+\epsilon_{\rm tail}\ll1$，经典支路另需 $\epsilon_{\rm tr,cl}^{\rm dyn}+\epsilon_{\rm tail}^{\rm cl}\ll1$ & $E_\ell^{(0)}=E_0+\ell\hbar\omega$ & 真实 $\mathcal E_\ell$、逐边 $\mathcal G_\ell,\Delta_\ell$ 多通道方程\\
双边平移 & 经典支路：$P_{\rm edge}\ll1$ 且 $\epsilon_{\rm edge}^{\rm dyn}=\Lambda_I\sqrt{P_{\rm edge}}\ll1$；量子支路：$\epsilon_{\rm edge,Q}^{\rm dyn}=\int\!\dd t\,|\upsilon(t)|\sqrt{\langle\Pi_{\min}(n+1)\rangle}\ll1$ & $\ell\in\mathbb Z$、$\hat b\hat b^\dagger=I$ & 有下界的单边通道链\\
低能 Dirac $\to$ Pauli & $\epsilon_{\rm NR}\ll1$，且正负能扇区混合受 $2m_{\mathrm e} c^2$ 能隙压低 & Pauli/Schr\"odinger 低能动力学 & 完整 Dirac 哈密顿量\\
散射 Born & $\epsilon_{\rm B,sc}=\|G_0^{(+)}V\|<1$（选定散射函数空间） & $T\simeq V$ 或有限 Born 阶数 & 精确 Lippmann--Schwinger $T=V+VG_0T$\\
二阶色散 & $\epsilon_{\rm disp}^{(\ge3)}(t)\ll1$ & 能量色散的二阶导数（Hessian）矩阵传播 & $\exp[-\ii E(\mathbf p)t/\hbar]$ 精确谱相位\\
近轴 & $\epsilon_{\rm par}\ll1$ 且 $\epsilon_{\rm par}^{(\ge4)}(z)\ll1$ & 二阶横向传播 & 精确相对论质量壳\\
开放系统 Born & $\epsilon_{\rm B,op}=g_B\tau_B\ll1$ & 二阶系统--环境核 & 精确联合 von Neumann 方程或 Nakajima--Zwanzig 方程\\
Markov & $\epsilon_{\rm M}=\tau_B/\tau_R\ll1$ & 时间局域记忆核 & Nakajima--Zwanzig\\
久期 & $\epsilon_{\rm sec}=(\Delta\omega_{\rm B}\tau_R)^{-1}\ll1$ & Bohr 频率分块的 GKSL & 非久期 Markov--Redfield 方程\\
\bottomrule
\end{longtable}


\section*{自由电子教学算例中的层级与记号}

动量连续态的振幅密度 $f_p(p)$ 满足 $\int|f_p|^2\dd p=1$，不能与正交离散通道的无量纲振幅 $c_\ell$ 混用。条件环境向量 $|\xi_r\rangle_X$ 一般未归一化，$\rho_{rs}=\langle\xi_s|\xi_r\rangle$；归一化环境态间的重叠记为 $\mathsf C_{rs}^X$。$\mathbf p_{\rm B}$ 是第一 Brillouin 区中的连续准动量代表，$\mathbf p_0$ 仍是入射束的物理中心动量。

格点在位失谐 $\delta_\ell=[E(p_\ell)-E_0]/\hbar-\ell\omega$ 与逐边失谐 $\Delta_\ell=\delta_{\ell+1}-\delta_\ell$ 都具有 $\mathrm{s^{-1}}$ 量纲，分别用于矩阵对角元与跃迁共振条件。单模投影矩阵元速率 $\mathcal G$ 与反厄米交换生成元的速率 $\upsilon$ 满足 $\upsilon=-\ii\mathcal G$，故 $g_{\mathrm Q}=\int\upsilon\dd t$ 无量纲。$r_Q=|g_{\mathrm Q}|^2$ 表示真空交换强度，$r_\Omega=|\Omega|/\omega_{\rm r}$ 表示驱动相对于反冲频率的比值。

无量纲周期密度统一写作 $\mathfrak n_e$，一个周期的平均为一；$n_e$ 继续用于单位体积内的电子数密度。$\theta_\zeta=\kappa_z\zeta$ 与探测面到达延迟 $\tau_{\rm d}$ 满足 $\theta_\zeta=-\omega\tau_{\rm d}$；漂移相位 $s_{\rm d}=\omega_{\rm r}L/v_0$ 无量纲。$\mathcal M_1=2|B_1|$ 是第一谐波调制幅度，允许大于一，不能无条件当作 Michelson 可见度。

正频场幅 $E_{\omega,z}$ 与实场峰值满足 $E_{\rm pk}^{\rm(real)}=2|E_{\omega,z}|$。参考层析中 $c_{01}=\rho_{01}$ 是未知电子相干，$b_R=|\beta_R|$ 是已标定参考耦合幅度；近场反演中的 $\beta_s$ 才是未知场积分。$\sigma_E$ 是独立高斯源宽与仪器宽的合成标准差，$\varepsilon_d=E_d-E_0$ 为读出能量偏差，$\Phi_N$ 是标准正态分布函数。$q_R$ 是由完整能谱方差构造的无量纲数据量，不是电荷。

热光输入的 $\bar n_{\rm th}$ 为平均光子数，$T_{\rm ph}$ 为模式温度；$\chi_e(\vartheta)$ 是电子阶数概率的特征函数，$\vartheta$ 为无量纲计数变量。$\lambda_\pm=r_Q(\bar n_{\rm th}+\tfrac12\mp\tfrac12)$ 分别对应增益和损失的 Poisson 表示参数，不是波长。四阶阶数累积量记为 $\mathcal C_4$。

高斯波包的 $C_p$ 与 $C_\zeta$ 分别表示动量和位置表象中的无量纲二次相位，不能仅凭同名“啁啾”将其互换。$a_0$ 是位置高斯指数的复参数，单位 $\mathrm{m^{-2}}$；$t_{\rm dif}$ 是自由展宽时间，$u_t=t/t_{\rm dif}$ 无量纲；$t_f,L_f$ 分别为焦点时间与传播距离。有限包络 $\eta_e(z)$ 按 $\int\eta_e\dd z=1$ 归一化，$Z_e=\int\eta_e\mathfrak n_e\dd z$ 修正完整线密度的归一化。双脉冲的 $\tau_p,D_p,\phi_p$ 分别为单脉冲持续时间、中心间隔及相对作用相位。

\section*{散射、传播与有限记忆噪声的记号}

推力变量 $T_{\rm thr}$、亏损 $\tau_T=1-T_{\rm thr}$、累计率 $\Sigma_T$、辐射权重 $R_T$ 和对数 $\mathcal L_T=\ln(1/\tau)$ 均无量纲。这里的 $\mathcal L_T$ 与表示 Talbot 长度的 $L_T$ 不同；定义见\eqnrefs{eq:qcd-thrust-definition,eq:qcd-thrust-radiator}。

中微子物质传播中 $V_{\rm CC}=\sqrt2G_F^{(E)}(\hbar c)^3n_e$ 为带电流势能；$n_e$ 的单位为 $\mathrm{m^{-3}}$。$\delta_\nu=\Delta m^2c^4/(2E_\nu)$ 为真空能量劈裂，$D_\nu$ 为物质中的瞬时本征能量间隔，均使用能量单位。$\theta_m$ 是物质混合角，$\epsilon_{\rm ad}=\hbar c|\theta_m'|/D_\nu$ 是本节的绝热判据，$r_n$ 是指数密度剖面的长度尺度；定义见\eqnrefs{eq:sm-matter-two-flavor,eq:sm-matter-adiabatic-condition}。PMNS 矩阵 $U_{\alpha i}$ 的行是带电轻子味，列是中微子质量指标，$W^+$ 带电流使用其复共轭。
强子动量份额 $\mathcal X_q,\mathcal X_g$ 是 PDF 的无量纲一阶动量矩，定义与演化见\eqnrefs{eq:qcd-momentum-fractions,eq:qcd-momentum-solution}；$\mu_{\rm F0}$ 是初始因子化动量尺度。DIS 的卷积归一化为 $\mathfrak F_2=F_2/x_B$、$\mathfrak F_L=F_L/x_B$，有限靶质量的物理纵向组合为 $F_L^{\rm phys}$。Drell--Yan 快度 $Y$ 与质量比 $\tau_{\rm DY}$ 均无量纲，定义见\eqnrefs{eq:dy-observables}；夸克分数电荷统一记为 $Q_f$。

中微子反冲例题使用入射能量 $E_\nu$、电子静止能量 $E_{e0}=m_{\mathrm e} c^2$ 和反冲动能 $T_e$。手征系数 $g_L^{(a)},g_R^{(a)}$ 为无量纲量；$C_{\nu e}$ 的单位为 $\mathrm{m^2/J}$。介子衰变常数 $f_P$、静止能量 $M_P,M_\ell$ 的单位均为能量，带电 pion 约定满足 $f_\pi=\sqrt2F_\pi$。相关矩阵元与率公式见\eqnrefs{eq:sm-pseudoscalar-axial-definition,eq:sm-pseudoscalar-width}。

频率噪声均方根 $\sigma_\varphi$ 的单位为 $\mathrm{s^{-1}}$，相关时间 $\tau_\varphi$ 的单位为秒。$C_\varphi$ 与双边谱 $S_\varphi$ 的单位分别为 $\mathrm{s^{-2}}$ 与 $\mathrm{s^{-1}}$。$\chi_\varphi$ 是滤波后相位方差的一半；无控制时与前文 $\Xi_\varphi$ 相同。阶差 $d_{\ell m}=\ell-m$ 和控制函数 $y_\varphi$ 无量纲，滤波振幅 $F_y$ 的单位为秒。由\eqnrefs{eq:electron-noise-filter-frequency} 可逐项检查指数无量纲。
